\chapter{Coronary Angiogram}

\section*{What Is This Test?}
A coronary angiogram, also called invasive coronary angiography or cardiac catheterization with coronary arteriography, is a hospital-based procedure that uses X-rays and iodinated contrast to show the coronary arteries. A thin catheter is introduced through an artery, usually in the wrist or groin, and advanced to the heart. Contrast is injected while X-ray images are recorded to identify narrowing, blockage, abnormal anatomy, or collateral blood flow. If a treatable blockage is found, angioplasty and stent placement may sometimes be performed during the same procedure, with separate consent and clinical assessment.

\section*{Alternative Names}
Coronary Angiography; Coronary Arteriography; Cardiac Catheterization; Heart Catheterization; Cardiac Cath; Invasive Coronary Angiogram; Diagnostic Coronary Angiogram

\section*{Principle}
Iodinated contrast makes the blood column visible under fluoroscopy. Sequential images show the course and lumen of the coronary arteries and help the cardiologist assess the location and severity of stenoses or occlusions. Pressure measurements, left-ventricular assessment, intravascular ultrasound, optical coherence tomography, or fractional flow reserve may be added when anatomy or the functional importance of a narrowing is uncertain. The procedure provides an anatomical and, when indicated, physiological assessment; it is not the same as a coronary CT angiogram.

\section*{Why This Test Is Ordered}
\begin{itemize}
  \item Suspected or confirmed acute coronary syndrome, including selected patients with heart attack or unstable angina
  \item Persistent or high-risk symptoms when noninvasive testing is abnormal, inconclusive, or not suitable
  \item Assessment before selected valve, congenital-heart, or other cardiac operations when coronary anatomy will affect management
  \item Planning revascularization with percutaneous coronary intervention (PCI) or coronary artery bypass grafting (CABG)
  \item Evaluation of selected unexplained ventricular dysfunction, malignant arrhythmia, or cardiac arrest when coronary disease is a possible cause
\end{itemize}

\section*{Primary vs Secondary Test}
Invasive coronary angiography is a definitive anatomical test used when the expected information will change management or when urgent assessment is needed. It is not a routine screening test for people without symptoms. Depending on clinical risk, alternatives may include ECG, echocardiography, exercise stress testing, stress imaging, or coronary CT angiography. A noninvasive test is often preferred for stable, lower-risk presentations when it can answer the clinical question safely.

\section*{How This Test Is Performed}
The patient receives local anesthetic and usually light sedation while remaining responsive. After sterile preparation, a sheath is placed in the radial artery at the wrist or the femoral artery at the groin. The catheter is guided under X-ray to the coronary arteries, contrast is injected, and images are obtained from multiple angles. The catheter and sheath are removed afterward, and pressure or a closure device is used to control bleeding. The procedure commonly takes about 30--60 minutes, though complex or combined procedures take longer.

\section*{What Preparation Is Needed}
Follow the cath-laboratory instructions about fasting and medications. Tell the team about kidney disease, diabetes, pregnancy or possible pregnancy, prior contrast reactions, bleeding disorders, and every medicine or supplement, especially anticoagulants and diabetes medicines. The team may check kidney function, blood count, electrolytes, coagulation tests, and an ECG. Arrange transportation and follow instructions about post-procedure observation; same-day discharge is possible after uncomplicated elective procedures, but emergency or interventional cases may require admission.

\section*{Contraindications and Precautions}
There is no absolute contraindication when urgent, life-saving information is required, but the risks and alternatives must be assessed. Important precautions include severe kidney impairment, previous severe contrast reaction, active bleeding or uncontrolled coagulopathy, severe anemia, infection at the access site, and pregnancy. Contrast exposure, hydration, medication adjustment, and access-site choice should be individualized. An apparent narrowing on angiography may not be flow-limiting, so pressure or intravascular assessment may be needed before deciding on intervention.

\section*{Risks}
Common effects include bruising, soreness, or a small lump at the access site. Less common but serious complications include bleeding requiring treatment, artery damage, contrast-associated kidney injury, allergic reaction, abnormal heart rhythm, blood clot, heart attack, stroke, infection, or death. The risk is higher in emergency procedures and in people with severe illness or complex vascular disease. Any planned stent or angioplasty adds procedure-specific risks and requires antiplatelet treatment planning.

\section*{What Happens After the Test}
The access site and circulation in the affected arm or leg are monitored, and fluids may be given when appropriate. The cardiologist discusses whether the arteries are normal, narrowed, or blocked and whether medication, PCI/stenting, CABG, or further testing is appropriate. Keep the access site clean and follow lifting and driving restrictions from the treating center. Seek urgent care for uncontrolled bleeding, increasing swelling, a cold or numb limb, severe chest pain, shortness of breath, fainting, fever, or a new neurologic symptom.

\section*{Understanding Results}

\subsection*{Normal or Nonobstructive}
The coronary arteries show no significant narrowing or only mild disease that does not materially restrict blood flow. This does not exclude coronary spasm, microvascular angina, early plaque, or non-cardiac causes of symptoms.

\subsection*{Abnormal}
The angiogram shows a focal or diffuse narrowing, complete occlusion, abnormal coronary origin, dissection, collateral vessels, or another structural finding. Severity is interpreted by the cardiologist in relation to symptoms, vessel location, plaque characteristics, pressure or imaging measurements, heart function, and overall risk. A narrowing does not automatically mean that a stent is needed; medical therapy or bypass surgery may be more appropriate in some patterns of disease.

\section*{Follow-up Tests}
\begin{itemize}
  \item \textbf{Fractional flow reserve or instantaneous wave-free ratio:} Measures whether an intermediate narrowing limits blood flow.
  \item \textbf{Intravascular ultrasound or optical coherence tomography:} Provides detailed images of the vessel wall and stent when anatomy is uncertain or intervention is planned.
  \item \textbf{Echocardiogram:} Assesses ventricular function, valves, and wall motion.
  \item \textbf{Kidney function and electrolyte tests:} Used when clinically indicated after contrast exposure, especially in patients at increased risk of kidney injury.
  \item \textbf{Cardiac rehabilitation and risk-factor assessment:} Includes blood pressure, lipid profile, glucose or HbA1c, smoking status, medication review, and individualized exercise planning.
\end{itemize}

\section*{References}
\begin{enumerate}
  \item American Heart Association. Cardiac Catheterization and Coronary Angiogram. Updated 2024.
  \item American Heart Association. Coronary Angiogram. Updated 2024.
  \item National Health Service. Angiography. Updated 2023.
\end{enumerate}

\clearpage
